\chapter{Conclusioni}
\label{cap:29-conclusioni}

Il \cref{cap:01-introduzione} ha aperto con una domanda: qual \`e la piattaforma
FaaS pi\`u piccola che sia ancora una piattaforma FaaS, e che cosa diventa
misurabile una volta ridotta a quella dimensione? Questo capitolo risponde alle
quattro domande di ricerca che ne discendevano.

\section{Le risposte}

\subsection{RQ1 --- Modularit\`a a costo zero sul percorso critico}

\textbf{S\`i, ed \`e verificabile.} La selezione modulare a tempo di build,
combinata con implementazioni no-op per ciascun punto di estensione, produce un
percorso di invocazione privo di rami condizionali sulla configurazione: la
catena di ripiego \emph{coda sincrona $\to$ coda asincrona $\to$ dispatch
diretto} non testa la presenza di un modulo, esegue il default del punto di
estensione precedente. La regola ArchUnit R3 rende la propriet\`a accertata dalla
build anzich\'e affermata.

Il costo esiste e \`e stato quantificato: due invocazioni della suite di test per
ogni verifica completa, e una copertura che riguarda i soli estremi dello spazio
delle configurazioni. La modularit\`a a costo zero sul percorso critico non \`e a
costo zero sulla catena di integrazione.

Le proporzioni, misurate di nuovo, non sono pi\`u quelle che questo capitolo
riportava: i dieci moduli sommano 6\,299 righe non vuote di produzione, contro
le 17\,334 del control plane e delle sue tre librerie obbligatorie
(\cref{cap:01-introduzione}). La riducibilit\`a della piattaforma non \`e quindi
la maggioranza del codice che si pu\`o non compilare, ma l'insieme delle
politiche che si possono non compilare; \`e un'affermazione pi\`u stretta e
verificata, non una proporzione.

\subsection{RQ2 --- Da quale segnale decidere}

\textbf{Dal tempo di sojourn, non da quello di servizio.} La risposta \`e
stabilita su due piani.

Analiticamente: il tempo di servizio cresce monotonamente con il limite di
concorrenza su qualunque risorsa condivisa, quindi una regola che riduce il
limite quando la latenza cresce non ha ottimo interno e cammina fino al proprio
pavimento. Il tempo di sojourn --- attesa pi\`u servizio --- l'ottimo interno ce
l'ha, perch\'e l'attesa decresce dove il servizio cresce.

Sperimentalmente: sotto carico l'attesa \`e stata misurata a 37--43\,ms contro un
servizio di circa 5\,ms, e un controllore pu\`o quindi sedere comodamente dentro
un SLO di servizio di 10\,ms mentre i suoi chiamanti attendono ottanta.

La realizzazione ha richiesto tre tentativi, e i due falliti sono contributi
quanto quello riuscito. Il gradiente sulla latenza end-to-end \`e instabile per
retroazione positiva. La ricerca del minimo per passi successivi \`e ben posta in
teoria e fallisce perch\'e non pu\`o distinguere ``la mia mossa ha aiutato'' da
``il carico \`e calato'' --- 36\,181 richieste rifiutate contro 32. Il progetto
che funziona \emph{calcola anzich\'e confrontare}: guardia di ammissione, legge di
Little su un tasso predetto, orizzonte di drenaggio proporzionale allo
scostamento dalla promessa. Sotto arrivi aperti ha servito il 3{,}2\% di
richieste in pi\`u, ne ha rifiutate il 15{,}7\% in meno e ha ridotto il p95
end-to-end del 10{,}6\% --- con \emph{meno} concorrenza media.

Il risultato \`e di una sola esecuzione per modo, e il
\cref{cap:27-discussione} lo classifica come suggerito e non stabilito.

\subsection{RQ3 --- Concorrenza come risorsa condivisa}

\textbf{Separare la stima del fabbisogno dall'allocazione funziona; la sua
superiorit\`a non \`e stata dimostrata.}

Il modo \code{BUDGETED} spezza la domanda in due --- quanto serve a una funzione
per il proprio SLO, quanto la piattaforma pu\`o dare --- e risolve la seconda con
equit\`a max-min pesata. Misurato, mantiene il totale sotto il budget e sposta
capacit\`a fra le funzioni al muoversi del carico, cosa che due controllori
indipendenti non possono fare per costruzione. Ha concesso alla funzione che non
ne aveva bisogno met\`a della concorrenza a parit\`a di throughput.

Due riserve. Il confronto \`e stato eseguito su un impianto a ciclo chiuso che,
come RQ4 ha poi mostrato, non poteva distinguere i segnali dei due controllori.
E i pesi --- l'argomento principale per cui il meccanismo esiste --- non sono mai
stati esercitati con valori diversi.

La campagna di co-tenancy ha per\`o prodotto il risultato pi\`u solido del lavoro,
e lo ha prodotto in negativo: il cross-talk fra funzioni co-residenti \`e reale
($-28\%$ di throughput sulla funzione stabile) e \textbf{non} \`e contesa di CPU
(vincolandole agli stessi quattro core, lo 0{,}6\% di differenza). La premessa
del modo a budget --- che esista una risorsa condivisa da dividere --- \`e
confermata; la sua identit\`a non \`e la CPU.

\subsection{RQ4 --- Il costo reale di un control plane}

\textbf{Meno di un core e circa 200\,MB di heap per 300 richieste al secondo, e
la scelta della catena di build conta pi\`u della scelta del framework.}

Tre risultati.

Per una JVM Spring, lo heap --- i dati --- \`e il 12\% della memoria residente. Il
resto \`e il runtime che descrive se stesso. Ottimizzare le allocazioni agisce
sull'ottavo del problema.

Il collector seriale di GraalVM Community spende il 31{,}8\% del tempo di parete
a raccogliere sotto carico sostenuto, con pause fino a 1{,}77\,s, e dimezza il
throughput. G1 lo elimina: 98{,}5\% del throughput della JVM, coda a 123\,ms,
avvio in 0{,}07\,s. Compilato a \code{-O3}, il control plane nativo supera la
JVM --- 9\,251 contro 8\,205 richieste al secondo, con un avvio ventitr\'e volte
pi\`u rapido.

Il terzo risultato \`e metodologico e vale oltre questo progetto: \emph{misure di
memoria senza un limite del container descrivono appetito, non fabbisogno}. Con
un limite a 1\,GB e uno heap esplicito, la stessa build consegna 3\,089
richieste al secondo in un terzo della memoria che l'esecuzione senza vincoli
aveva occupato.

\section{La tesi, rivista}

La tesi del \cref{cap:01-introduzione} era che un insieme selezionato di rinunce
produce un sistema pi\`u utile come strumento di misura di quanto lo sarebbe la
sua versione completa. A valle dei risultati, il bilancio \`e disomogeneo e vale
la pena dichiararlo.

Il \textbf{processo singolo} ha pagato pi\`u di quanto promesso. Ha reso
attribuibile una differenza di pochi millisecondi di attesa a una politica di
ammissione --- che \`e l'intero contenuto del \cref{cap:25-valutazione-concorrenza}
--- e si \`e rivelato, per esclusione, anche il sospettato principale del
cross-talk fra funzioni co-residenti. Il vincolo ha prodotto lo strumento e il
fenomeno.

Lo \textbf{stato in memoria} ha pagato in parte, a un costo pi\`u alto del
previsto: tre orizzonti di ritenzione e un intero apparato di dimenticanza per
evitare che uno store che non persiste diventi uno store che non dimentica.

L'\textbf{assenza di autenticazione} non ha pagato. Nessun risultato di questo
documento vi dipende, ed \`e la sola rinuncia che impedisca l'uso fuori da un
laboratorio.

\section{Il contributo che non era previsto}

Il contributo di cui vado pi\`u convinto non compare fra i sei elencati nel
\cref{cap:01-introduzione}, ed \`e la pratica di conservare i risultati negativi
insieme alle affermazioni che hanno confutato.

Tre dei risultati principali sono negativi. Ciascuno ha richiesto
un'implementazione completa e una campagna di misura, e ciascuno ha eliminato una
famiglia intera di soluzioni apparentemente ragionevoli: il gradiente sul tempo
di servizio, la ricerca del minimo per confronto fra istanti, il carico a ciclo
chiuso per valutare politiche di ammissione. A questi si aggiungono un errore di
calibrazione che ha fabbricato 86\,302 rifiuti, un test invalidato da un flag
inefficace, e un confondente --- due funzioni omonime con runtime diversi ---
che ha invalidato analisi precedenti.

Nessuno di questi \`e derivabile dalla lettura del codice finale, che per
costruzione contiene solo ci\`o che ha funzionato. Il costo di non registrarli \`e
che la stessa famiglia di soluzioni viene ritentata, e il segnale che le esclude
va riprodotto da capo.

\section{Chiusura}

\nf{} \`e piccolo di proposito. Il suo percorso critico attraversa un solo
processo, e le sue politiche --- code, autoscaling, controllo della concorrenza,
offload --- si possono non compilare. Non compete con le piattaforme mature su completezza e
non aspira a farlo.

Ci\`o che ha prodotto \`e una risposta precisa a una domanda che le piattaforme
complete rendono difficile porre: da quale grandezza misurata debba decidere un
controllore di ammissione, e che cosa succede quando decide dalla grandezza
sbagliata. La risposta --- il tempo di sojourn, non quello di servizio --- \`e
banale una volta enunciata, ed \`e costata due progetti falliti e la scoperta che
il generatore di carico stava impedendo alla domanda di essere posta.

\`E il tipo di risultato che una piattaforma abbastanza piccola da essere letta
per intero consente di ottenere, e che una piattaforma abbastanza grande da
essere utile in produzione avrebbe nascosto sotto tre livelli di proxy.
